\begin{proposition}[6.12.9]
\label{IV.2.s6.12.9-fr}
Sea $X$ un preesquema localmente noetheriano tal que $\operatorname{Reg}(X)$
es abierto; entonces, para todo $n\geq0$, el conjunto $U_{R_n}(X)$ de los
puntos $x\in X$ en los que $X$ satisface $(R_n)$ es abierto.
\end{proposition}

\begin{proof}
En efecto (5.8.2), $U_{R_n}(X)$ es la unión de $\operatorname{Reg}(X)$ y
del conjunto de los $x\in\operatorname{Sing}(X)$ tales que los puntos
genéricos $z_i$ de las componentes irreducibles de $\operatorname{Sing}(X)$
que contienen a $x$ satisfacen la relación
$\dim(\mathcal O_{X,z_i})\geq n+1$ (puesto que
$z_i\notin\operatorname{Reg}(X)$); en otras palabras, estos puntos
$x\in\operatorname{Sing}(X)$ son aquellos para los que se tiene
$\operatorname{codim}_x(\operatorname{Sing}(X),X)\geq n+1$ (5.1.2); como
la función $x\leadsto\operatorname{codim}_x(\operatorname{Sing}(X),X)$ es
semicontinua inferiormente (\textbf 0, 14.2.6), el conjunto de estos puntos es
abierto en $\operatorname{Sing}(X)$, lo que termina la demostración.
\end{proof}
